\chapter{Statistical foundations and relation to prior work}
\label{ch:lineage}
\chapterpromise{The established literature on exponential families, product partitions, exact Bayesian changepoint recursion, irregular designs, multiple sequences, latent groups, and numerical approximation.}

\section{Exponential families and conjugate analysis}
For a regular exponential family, likelihood contributions can be represented through sufficient
statistics, a natural parameter, and a log-partition function. Conjugate priors for exponential
families preserve the natural form of the posterior and express integrated likelihoods through
normalizing-constant ratios \citep{diaconis1979conjugate,bernardo1994bayesian}. This is established
Bayesian theory, not a new property of \BayesBreak. Its role here is to provide exact segment
marginal likelihoods and posterior moments for every candidate interval when the required
normalizers are finite. Exposure in Poisson sampling and trial counts in Binomial sampling enter the
likelihood through family-specific sufficient summaries; they are not generic likelihood weights.

\section{Product-partition models and exact Bayesian changepoint inference}
Product-partition models assign a probability to a partition through products of segment cohesion
terms and combine this prior with segmentwise integrated likelihoods
\citep{barry1992ppm,barry1993bayesCP}. Exact Bayesian multiple-changepoint recursions exploit the
same factorization to sum over possible previous changepoints and compute posterior distributions
without enumerating every partition explicitly \citep{fearnhead2006exact}. Bayesian
piecewise-constant regression further shows how partition averaging yields posterior regression
curves, changepoint probabilities, and marginal likelihoods \citep{hutter2006bpcr}.

Accordingly, this book does not claim to introduce product partitions, conjugate segment
integration, or dynamic programming for Bayesian changepoint models. \BayesBreak uses these
established results as the foundation for a broader model specification covering multiple
observation families, design-dependent priors, multiple aligned sequences, and group structure.
The sum-product recursion for posterior marginals and the max-sum recursion for a joint MAP
partition are stated separately because they solve different algebraic problems even when they use
the same segment scores.

\section{Frequentist segmentation and alternative Bayesian computation}
Frequentist methods optimize additive segment costs or multiscale criteria. Segment-neighborhood
and optimal-partitioning algorithms provide exact minimizers for fixed segment counts or penalties
\citep{auger1989segment,jackson2005optpart}; pruned dynamic programming and PELT reduce practical
cost under additional conditions \citep{rigaill2010pruned,killick2012pelt}. Total-variation and
fused-lasso estimators use convex regularization, including group penalties for multiple sequences
\citep{rudin1992rof,tibshirani2005fused,bleakley2011groupfused}. SMUCE and wild binary segmentation
provide multiscale or recursive detection procedures with different inferential targets
\citep{frick2014smuce,fryzlewicz2014wbs}. Circular binary segmentation is widely used for array-CGH
analysis \citep{olshen2004cbs}. These procedures are relevant comparators but do not, without
additional construction, return the same posterior distribution over partitions.

When segment parameters cannot be marginalized or the dependence structure is not amenable to a
finite recursion, reversible-jump Markov chain Monte Carlo and other trans-dimensional sampling
methods provide a general Bayesian alternative \citep{green1995rjMCMC,denison1998bayesian}. Bayesian
online changepoint detection instead filters a run-length distribution as data arrive and solves a
different online problem \citep{adamsmackay2007bocpd}. \BayesBreak is an offline finite-partition
method.

\section{Irregular designs and design-dependent partition priors}
Product-partition models have been generalized so that cohesion depends on covariates or other
predictors \citep{park2010gppm,muller2011ppmx}. In changepoint problems, irregular observation
coordinates affect the prior plausibility of segment lengths and boundary locations. This book
therefore distinguishes a segment-cohesion factor $c_u(i,j)$ from an interior-boundary hazard
$h_u(j)$. The distinction prevents two different prior constructions---a prior on segment length and
a local hazard of a changepoint---from being described as the same object. It also keeps design
spacing separate from exposure or trial information in the observation likelihood.

\section{Hierarchical and multiple-sequence changepoint models}
Hierarchical Bayesian changepoint models have a long history
\citep{carlin1992hierBayesCP}. Methods for simultaneous changepoints in multiple data sequences use
hierarchical priors to borrow information about the propensity for changes across sequences
\citep{fan2017basic}. Joint random partition models provide more flexible dependence among
partitions across several ordered processes \citep{quinlan2024jrpm}. Other work permits dependence
between successive segments or richer within-segment regression structures
\citep{fearnhead2011dependence,ruggieri2014bcpvs}.

The exact multiple-sequence calculation developed here is deliberately narrower: aligned sequences
share the same candidate partition, are conditionally independent given that partition, and may
have sequence-specific segment parameters. Their segment marginal likelihoods then multiply. This
is an exact hierarchical model under those assumptions, not a claim to encompass every form of
partial synchrony or dependence among sequence-specific changepoints. Known groups repeat the
calculation within each group.

\section{Latent groups and optimization}
Model-based clustering of longitudinal or ordered data with common structural changes is relevant
prior work for clustering sequences through changepoint patterns
\citep{corradin2026commonstructural}. Finite-mixture identifiability results concern normalized
component distributions \citep{teicher1963identifiability}, and the EM algorithm concerns
likelihoods or objectives for which the auxiliary-function construction is valid
\citep{dempster1977em}.

The latent-group calculation in this book is stated more narrowly. It defines a finite objective
\[
\mathcal F(\pi,\tau)=\sum_s\log\!\left\{\sum_g\pi_gS_s(\tau_g)\right\},
\]
where $S_s(\tau_g)$ is a positive sequence-template score. A Jensen minorizer yields auxiliary
allocation weights; exact maximization of the minorizer over template partitions uses a
responsibility-weighted max-sum recursion. The proof establishes monotone ascent of this stated
criterion. It does not establish that $S_s(\tau_g)$ is a normalized component density, identify a
finite mixture, or average over uncertainty in the templates.

\section{Nonconjugate segment models}
Laplace approximation, variational bounds, P\'olya--Gamma augmentation, expectation propagation,
and numerical quadrature are established tools for nonconjugate Bayesian calculations
\citep{tierney1986laplace,jaakkola2000logisticvb,polson2013pg,minka2001ep}. \BayesBreak uses them, when
appropriate, to approximate individual segment marginal likelihoods. The general theorem in
Chapter~\ref{ch:nonconj} is conditional: if every admissible segment log marginal likelihood is
within $\varepsilon$ of a reference value, then partition marginal likelihoods and posterior odds
obey explicit bounds. The theorem does not supply a method-specific value of $\varepsilon$.

\section{Position of the present work}
The contribution is a generalized hierarchical Bayesian segmentation formulation built from
established components whose assumptions are kept explicit. It combines:
\begin{enumerate}
\item exact segment integration for a class of conjugate exponential-family observation models;
\item sum-product posterior computation and separate max-sum MAP optimization over ordered
partitions;
\item priors for irregular designs that distinguish segment cohesion from boundary hazard;
\item exact common-changepoint pooling for aligned conditionally independent sequences and known
groups;
\item a finite latent-group template criterion with a proved monotone alternating algorithm;
\item conditional propagation of numerical segment-score error; and
\item family-specific posterior prediction and clearly defined changepoint metrics.
\end{enumerate}
The novelty claim is attached to this combined generalized method and its derivations, not to the
prior existence of dynamic programming, exponential-family conjugacy, product partitions, or EM.

\begin{figure}[H]
\centering
\resizebox{0.96\textwidth}{!}{
\begin{tikzpicture}[node distance=8mm and 10mm]
\node[successbox,text width=28mm] (block) {Block integration\\established};
\node[successbox,right=of block,text width=28mm] (dp) {Sum-product identities\\established};
\node[successbox,right=of dp,text width=28mm] (map) {Max-sum correctness\\established};
\node[secondarybox,below=of dp,text width=28mm] (prior) {Cohesion + hazard\\repaired};
\node[successbox,below=of map,text width=28mm] (shared) {Finite-candidate consistency\\established under conditions};
\node[secondarybox,below=of block,text width=28mm] (score) {Score-template monotonicity\\repaired};
\node[primarybox,below=of prior,text width=28mm] (approx) {Segment-score error propagation\\conditional};
\draw[arrPrimary] (block)--(dp);\draw[arrPrimary] (dp)--(map);
\draw[arrPrimary] (prior)--(dp);\draw[arrPrimary] (block)--(shared);\draw[arrPrimary] (dp)--(shared);
\draw[arrPrimary] (block)--(score);\draw[arrPrimary] (block)--(approx);\draw[arrPrimary] (dp)--(approx);
\end{tikzpicture}
}
\caption{Logical dependence of the principal results. Exact partition calculations require finite
segment marginal likelihoods and factorized prior contributions. The nonconjugate result additionally
requires a uniform segmentwise error bound.}
\label{fig:theorem-dependency}
\end{figure}

\section{Recommended reading sequence}
A reader new to the subject should begin with exponential-family conjugacy, product-partition
models, and exact Bayesian multiple-changepoint recursion; then proceed to Bayesian
piecewise-constant regression, design-dependent product partitions, and hierarchical
multiple-sequence methods. The literature annotations in Appendix~\ref{app:annotated-lit} state the
specific role of every reference retained in the book. Numerical implementation additionally
requires stable log-sum-exp evaluation and the approximation methods used for nonconjugate segment
models \citep{blanchard2021lse}.
